\chapter{Banking decisions, review effort and the value of assistance}
\label{la:economics}

\section{Value belongs to a decision and an information set}

The value of an agent is the change it makes to a feasible decision after
including its information, review and operating costs. Producing a longer
report, drawing more edges or executing a notebook faster can help that
decision, but none is the economic target itself. The banking application must
specify the action, baseline workflow, affected population, financial horizon
and loss before comparing systems.

Let $a$ range over feasible actions and let $V(a,\theta)$ be their monetary
value under a causal state $\theta$. An information set $\mathcal I$ might
contain observed data, authenticated design records and unresolved alternatives.
If a justified probability law for $\theta$ is available, an expected-value
decision solves
\[
 a^*(\mathcal I)\in\arg\max_a
       \mathbb E\{V(a,\theta)\mid\mathcal I\}.
\]
This expression is a decision rule conditional on a probability model. A
language model's self-reported certainty, the frequency of repeated answers,
or a fraction of retrieved documents is not automatically that probability.
The distinction matters most when a small change in the assumed probability
changes a lending, servicing or funding action.

If the evidence instead supports a set $\Theta(\mathcal I)$ without probabilities,
the bank needs an explicit ambiguity policy. For example, a robust rule can
maximize $\inf_{\theta\in\Theta(\mathcal I)}V(a,\theta)$. Expected value and this
robust criterion answer different decision questions. The software must
preserve the information needed for either and require the policy to be named.
It must not manufacture a distribution by assigning equal weight to whichever
graphs an algorithm happened to enumerate.

\section{An exact example of useful and costly information}

Consider a stipulated one-customer intervention whose incremental gross return
is either CNY~0 or CNY~400, with intervention cost CNY~300. Suppose the two
states have equal prior probability. These are project assumptions for an
arithmetic example, not bank estimates, elicited probabilities or measured
model performance. Acting has expected net return $200-300=-100$, so the
unassisted expected-value decision is to do nothing, with value zero.

Perfect information would permit action only in the CNY~400 state. Its expected
gross value relative to doing nothing is therefore
\[
 \tfrac12\max(0,0-300)+\tfrac12\max(0,400-300)=50.
\]
For this action set and loss, no information service can justify an expected
payment above CNY~50 per decision solely through improved action selection.
The bound would change if the service had other validated benefits or changed
the feasible actions; neither is assumed here.

Now stipulate a binary signal that reports ``positive'' with probability $0.9$
in the high-return state and $0.1$ in the low-return state. Each signal outcome
has marginal probability $0.5$. Bayes' rule gives posterior high-state
probabilities $0.9$ after a positive signal and $0.1$ after a negative signal.
The conditional net returns from acting are respectively
\[
 0.9(400)-300=60,\qquad 0.1(400)-300=-260.
\]
The optimal rule acts only after a positive signal. Its gross expected value
is $0.5(60)=30$. An independent calculation keeps the loss from false positives
visible:
\[
 \underbrace{0.5(0.9)(100)}_{\text{true-positive gain}}
 -\underbrace{0.5(0.1)(300)}_{\text{false-positive loss}}=30.
\]
The information value is below the perfect-information bound, as it should be.

Suppose obtaining and checking the signal costs CNY~2 in computation and CNY~8
in review for every request, including negative ones. The incremental net value
is then CNY~20. Under the additional simplifying assumptions of identical
expected value, linear costs and no discounting, a CNY~50,000 fixed development
cost breaks even at $50{,}000/20=2{,}500$ decisions. These figures are reproduced
with exact rational arithmetic. They are a transparent planning construction,
not a price quotation or an ROI estimate for the implemented prototype.

The threshold for acting is a high-state probability above $300/400=0.75$.
An assistant that says ``90\% confident'' can cross that threshold without
having a calibrated probability. Review must establish the meaning and
reliability of the signal. A syntactically valid report, a correct posterior
formula and an accurate multiplication cannot compensate for an invented
likelihood.

\section{Refusal can preserve value, but it also consumes resources}

If all that is defensible is the set $\Theta=\{0,400\}$, robust net value from
acting is $-300$, while doing nothing guarantees zero incremental value.
The robust rule therefore declines the intervention. This conclusion differs
from a claim that the true effect is zero. It expresses how the bank chooses
under unresolved uncertainty and the stated loss.

A service that exposes the ambiguity can be valuable if the incumbent workflow
would act on a falsely precise estimate. The evaluation must include that
incumbent behavior explicitly. Against an already cautious baseline, the same
refusal may add little decision value while still costing analyst time. A
refusal is neither automatically a failure nor automatically a benefit.

The inspected observational workflow illustrates this problem through critic
declines under poor overlap~\citep{v12_la_2607_22443}. The denominator must include
the declined requests, their review cost and what happened next. Did a reviewer
obtain a better design, narrow the population with explicit authorization,
commission an experiment, or simply lose time? Each path has a different
economic implication. Changing to a selected population must also change the
target rather than make the original request appear successful.

\section{Four banking journeys require different evidence}

For a funding decision, assistance can assemble the rate action, balance units,
margin horizon and policy constraints. It can also flag a mismatch between
calendar covariance and an assumed independent sampling unit. Value arises
only if these corrections improve a decision or avert an error relative to the
baseline. A conservative interval wider than the bank's action threshold may
be scientifically defensible and economically inconclusive.

For sequential servicing, the central object is a policy adapted to information
available at each decision. The total bill includes state measurement,
assignment logging, intervention delivery and customer outcomes, as well as
model calls. An agent can help describe a dynamic regime, but a static ATE
translator cannot evaluate it by substituting a single treatment indicator.
The economic study must compare genuinely feasible policies under common
information and count unsuccessful or unavailable evaluations.

For discovery followed by intervention, a visually simpler graph can conceal
alternative effects. The useful information is whether additional evidence
changes the set of financially acceptable actions. More accurate recovery of
irrelevant edges may have little value for the selected decision. Conversely,
one mistaken orientation near the treatment or outcome can reverse a policy
recommendation. Graph accuracy and decision loss should both be reported, with
neither used as a proxy for the other.

For trustworthy observations, a language system can reconcile field names and
suggest ownership or inclusion questions. The bank must still establish actual
linkage, label and audit truth. Reviewing a source dictionary may prevent a
semantic error; it does not supply sampled identity labels or independent
calendars. The cost of obtaining those scientific inputs belongs in the
information budget, and their absence remains a legitimate stop condition for
an effect calculation.

\section{A complete cost model}

For request $i$, a prospective accounting identity is
\[
 C_i=C_{i,\mathrm{model}}+C_{i,\mathrm{retrieval}}+
     C_{i,\mathrm{numeric}}+C_{i,\mathrm{review}}+
     C_{i,\mathrm{data}}+C_{i,\mathrm{operation}}.
\]
The terms represent distinct resource uses, not independently measured values
in the current project. Model cost includes all calls and retries. Retrieval
includes acquisition and any permitted indexing. Numerical cost includes the
complete fit, uncertainty calculation and serialization. Review includes time
spent on wrong, refused and unsupported answers. Data cost includes audits and
additional information. Operations include storage, recovery and support.

A natural-workflow comparison measures the quality and cost of complete
workflows. A matched-budget comparison holds selected resources fixed to study
one contribution. The two can produce different rankings without contradiction.
An agent that improves answers by using more expert review may be valuable,
but the result is about the combined workflow. It does not isolate autonomous
reasoning quality.

The experimental unit must also match the cost and outcome calculation.
Repeated prompts about one mechanism are correlated observations within that
mechanism. Counting them as independent bank decisions understates uncertainty
and can exaggerate the amount of evidence. Error bars for incremental value
should use the prespecified mechanism or decision clusters and retain the
paired comparison wherever both workflows face the same problem.

\section{What can be concluded before a live evaluation}

The local controller demonstrates preservation of declarations, alternatives,
refusals and numerical set topology in selected authored examples. It makes
zero live model calls. Its measured interface behavior cannot reveal a
provider's accuracy, token use, price or review burden on unfamiliar questions.
Those fields remain unavailable rather than being filled with zeros.

Published benchmarks support the design of a prospective evaluation. They
cannot be combined across different model revisions, task distributions and
information budgets into a bank ROI estimate. CausalDS in particular reports
divergence between symbolic success, interval quality, abstention and cost,
which supports measuring them separately~\citep{v12_la_2607_08093}.

The next empirical step requires a specified backend and custody arrangement,
actual prices and budget, independent held-out mechanisms, adjudicators and
qualified numerical baselines. For bank adoption it additionally requires
feasible actions, authentic records and accepted loss thresholds. Until those
inputs exist, exact decision examples can check the accounting and expose
threshold errors. They cannot establish that assistance crosses an institutional
investment threshold.
